% =============================================================================
%  Capítulo 5 -- Propuesta del Índice de Neutralidad Sociolectal
% =============================================================================
\chapter{Propuesta del Índice de Neutralidad Sociolectal}
\label{cap:ins}

\section{Definición conceptual del índice}
\label{sec:ins-conceptual}

El \emph{Índice de Neutralidad Sociolectal} (\INS) busca operacionalizar una
intuición sencilla: un sistema de clasificación de texto es tanto más
\emph{neutral} respecto a la variación sociolectal cuanto más consistente sea
su comportamiento al ser confrontado con pares de textos que comparten
contenido informativo pero difieren en registro lingüístico.

El \INS no es una métrica de \emph{exactitud}, sino de \emph{consistencia
intra-par}. Un modelo puede tener exactitud baja y, sin embargo, \INS alto si
sus errores se distribuyen de manera independiente del registro; y un modelo
puede tener exactitud alta y, no obstante, \INS bajo si su comportamiento
correcto se concentra sistemáticamente en un registro y no en otro.

\section{Formalización preliminar}
\label{sec:ins-formal}

\begin{definicion}[Par controlado]
\label{def:par-controlado}
Un \emph{par controlado} es un par ordenado $(t, t')$ de textos tales que su
contenido informativo, evaluado bajo los criterios de la
Sección~\ref{sec:benchmark-diseno}, es operativamente equivalente, y que
difieren únicamente en su registro lingüístico.
\end{definicion}

Sea $\mathcal{B} = \{(t_i, t_i')\}_{i=1}^{N}$ un \emph{benchmark} de pares
controlados y sea $f$ un clasificador. Se proponen dos formulaciones
complementarias del \INS, según la información de salida disponible.

\subsection{Formulación basada en decisiones}

Cuando solo se dispone de la decisión $f(t) \in \mathcal{Y}$:

\begin{definicion}[\INS de decisión]
\label{def:ins-decision}
\[
  \mathrm{INS}_{\mathrm{dec}}(f, \mathcal{B}) \;=\;
  \frac{1}{N} \sum_{i=1}^{N} \mathbb{1}\!\left[ f(t_i) = f(t_i') \right].
\]
\end{definicion}

\noindent
Esta formulación corresponde a la proporción de pares en los que el modelo
emite la misma decisión para ambos miembros. Toma valores en $[0,1]$, con $1$
indicando consistencia perfecta de decisión.

\subsection{Formulación basada en probabilidades}

Cuando se dispone de la distribución posterior $p_f(\cdot \mid t)$ sobre las
clases:

\begin{definicion}[\INS probabilístico]
\label{def:ins-prob}
\[
  \mathrm{INS}_{\mathrm{prob}}(f, \mathcal{B}) \;=\;
  1 \;-\; \frac{1}{N} \sum_{i=1}^{N}
  d\!\left( p_f(\cdot \mid t_i),\; p_f(\cdot \mid t_i') \right),
\]
\noindent donde $d(\cdot,\cdot)$ es una medida de divergencia entre
distribuciones (por ejemplo, divergencia de Jensen--Shannon o distancia de
variación total), normalizada al intervalo $[0,1]$.
\end{definicion}

\noindent
La formulación probabilística captura cambios en la confianza del modelo que
la formulación de decisión ignora, y por tanto es preferible cuando se dispone
de los puntajes de probabilidad.

\subsection{Propiedades deseables}

\begin{proposicion}[Acotación]
$\mathrm{INS}_{\mathrm{dec}}(f, \mathcal{B}), \mathrm{INS}_{\mathrm{prob}}(f, \mathcal{B}) \in [0,1]$.
\end{proposicion}

\begin{proposicion}[Invariancia a permutaciones del \textit{benchmark}]
El valor de ambas formulaciones es invariante a cualquier reordenamiento de
los pares de $\mathcal{B}$.
\end{proposicion}

[Discutir y, según corresponda, demostrar otras propiedades: comportamiento
ante \emph{benchmarks} desbalanceados por tipo de variación, robustez a
ruido, etc.]

\subsection{Variantes desagregadas}

Si los pares de $\mathcal{B}$ están etiquetados por tipo de variación
lingüística $v \in V$, se define el \INS condicional como

\[
  \mathrm{INS}_{\mathrm{dec}}(f, \mathcal{B} \mid v)
  \;=\;
  \frac{1}{|\mathcal{B}_v|}
  \sum_{(t,t') \in \mathcal{B}_v}
  \mathbb{1}\!\left[ f(t) = f(t') \right],
\]
\noindent donde $\mathcal{B}_v$ es el subconjunto de pares cuyo tipo de
variación es $v$. Esta variante permite identificar qué tipos de variación
afectan más al modelo.

\section{Procedimiento de cálculo}
\label{sec:ins-procedimiento}

Esta sección operacionaliza las definiciones de la
Sección~\ref{sec:ins-formal} en un procedimiento ejecutable y reproducible.
Se describen las entradas requeridas, el algoritmo paso a paso, las
decisiones de implementación y una implementación de referencia.

\subsection{Entradas y supuestos}

El cálculo del \INS requiere dos entradas. La primera es un \emph{benchmark}
$\mathcal{B} = \{(t_i, t_i')\}_{i=1}^{N}$ de pares controlados
(Definición~\ref{def:par-controlado}), donde cada par está opcionalmente
etiquetado con su tipo de variación lingüística $v \in V$ para permitir la
desagregación. La segunda es un clasificador $f$ que expone, como mínimo, una
interfaz de decisión $f(t) \in \mathcal{Y}$ y, de manera opcional, la
distribución de probabilidad posterior $p_f(\cdot \mid t)$ sobre las clases.

El procedimiento presupone que las salidas del modelo son deterministas o han
sido vueltas reproducibles. Para modelos con componentes estocásticos ---en
particular, modelos de lenguaje grandes usados como clasificadores--- esto
exige fijar la temperatura de muestreo en cero o, cuando ello no sea posible,
promediar las salidas sobre $k$ ejecuciones independientes y documentar $k$.
Sin esta precaución, la inconsistencia atribuible al azar del muestreo se
confundiría con la inconsistencia sociolectal que el índice pretende medir.

\subsection{Algoritmo}

El cálculo procede en cinco pasos.

\begin{enumerate}
  \item \textbf{Verificación del \emph{benchmark}.} Se comprueba que
    $\mathcal{B}$ no esté vacío, que cada par esté completo y que las
    etiquetas de tipo de variación, si existen, pertenezcan al catálogo $V$
    definido en la Sección~\ref{sec:benchmark-diseno}. Un \emph{benchmark}
    vacío no produce un \INS de cero sino un valor \emph{indefinido}: el
    procedimiento debe fallar explícitamente en ese caso, no devolver un
    número que pueda malinterpretarse.
  \item \textbf{Obtención de salidas.} Para cada par $(t_i, t_i')$ se
    obtienen $f(t_i)$ y $f(t_i')$ y, si el modelo lo permite,
    $p_f(\cdot \mid t_i)$ y $p_f(\cdot \mid t_i')$. Esto implica exactamente
    $2N$ invocaciones al modelo; el costo del procedimiento es lineal en el
    tamaño del \emph{benchmark}.
  \item \textbf{Cálculo de la formulación de decisión.} Se computa
    $\mathrm{INS}_{\mathrm{dec}}$ según la Definición~\ref{def:ins-decision}:
    la proporción de pares con decisión coincidente.
  \item \textbf{Cálculo de la formulación probabilística.} Si se dispone de
    probabilidades, se computa $\mathrm{INS}_{\mathrm{prob}}$ según la
    Definición~\ref{def:ins-prob}, con la divergencia de Jensen--Shannon como
    medida $d$ (la elección se justifica más abajo).
  \item \textbf{Desagregación y reporte.} Se calcula el \INS condicional por
    tipo de variación, $\mathrm{INS}(f, \mathcal{B} \mid v)$, y se reporta el
    índice global junto con el perfil desagregado, el tamaño $N$ del
    \emph{benchmark}, el tamaño $N_v$ de cada celda y las métricas de
    desempeño del modelo sobre la tarea. La razón de este último requisito se
    desarrolla en la Sección~\ref{sec:ins-interpretacion}.
\end{enumerate}

\subsection{Decisiones de implementación}

\paragraph{Elección y normalización de la divergencia.} La
Definición~\ref{def:ins-prob} exige una divergencia $d$ acotada en $[0,1]$.
Se adopta la divergencia de Jensen--Shannon calculada con logaritmo en base 2,
que satisface esta cota de manera natural ---alcanza $0$ para distribuciones
idénticas y $1$ para distribuciones con soportes disjuntos--- sin requerir
normalización \emph{ad hoc}. Es además simétrica y está definida incluso
cuando alguna clase recibe probabilidad cero, pues la mezcla
$m = \tfrac{1}{2}(p + p')$ es estrictamente positiva dondequiera que alguna
de las dos distribuciones lo sea.

\paragraph{Degradación controlada.} Cuando el modelo no expone
probabilidades ---caso frecuente en servicios de clasificación accedidos como
caja negra--- el procedimiento se degrada a la formulación de decisión y el
reporte lo hace explícito. Comparar un $\mathrm{INS}_{\mathrm{prob}}$ de un
modelo con un $\mathrm{INS}_{\mathrm{dec}}$ de otro no es válido: son
cantidades distintas.

\paragraph{Separación entre evaluación y métrica.} La implementación separa
la obtención de salidas del modelo (paso 2) del cálculo del índice (pasos
3--5). Esta separación permite probar unitariamente las funciones de cálculo
con casos sintéticos de respuesta conocida y recalcular el índice sin repetir
las $2N$ invocaciones al modelo, que son la parte costosa del procedimiento.

\subsection{Implementación de referencia}

El Código~\ref{lst:ins-ref} presenta la implementación de referencia en
Python, sin más dependencia que \texttt{NumPy}. Las funciones de cálculo
fueron verificadas con pruebas unitarias que cubren los casos extremos:
consistencia perfecta ($\INS = 1$), inversión total ($\INS = 0$), cotas de la
divergencia y rechazo del \emph{benchmark} vacío.

\begin{lstlisting}[caption={Implementación de referencia del cálculo del \INS.},label={lst:ins-ref}]
import numpy as np

def divergencia_js(p, q, base=2.0):
    """Divergencia de Jensen-Shannon entre dos distribuciones
    discretas. Con base 2, el resultado queda acotado en [0, 1]."""
    p = np.asarray(p, dtype=float); q = np.asarray(q, dtype=float)
    p = p / p.sum(); q = q / q.sum()
    m = 0.5 * (p + q)
    def kl(a, b):
        mascara = a > 0   # m > 0 donde a > 0: el cociente es seguro
        return float(np.sum(a[mascara]
                     * (np.log(a[mascara] / b[mascara]) / np.log(base))))
    return 0.5 * kl(p, m) + 0.5 * kl(q, m)

def calcular_ins_dec(pares):
    """INS de decision. `pares` es una lista de tuplas (y, y_prima)
    con las decisiones del modelo para cada par controlado."""
    if not pares:
        raise ValueError("Benchmark vacio: el INS no esta definido.")
    return sum(int(y == yp) for y, yp in pares) / len(pares)

def calcular_ins_prob(pares, base=2.0):
    """INS probabilistico. `pares` es una lista de tuplas (p, p_prima)
    con las distribuciones posteriores del modelo para cada par."""
    if not pares:
        raise ValueError("Benchmark vacio: el INS no esta definido.")
    total = sum(divergencia_js(p, q, base) for p, q in pares)
    return 1.0 - total / len(pares)

def ins_por_variacion(registros):
    """INS de decision desagregado por tipo de variacion.
    `registros` es un iterable de tuplas (tipo, y, y_prima)."""
    grupos = {}
    for tipo, y, yp in registros:
        grupos.setdefault(tipo, []).append((y, yp))
    return {t: calcular_ins_dec(ps) for t, ps in grupos.items()}
\end{lstlisting}

\subsection{Ejemplo ilustrativo}

Para fijar ideas, considérese un \emph{benchmark} ficticio de cuatro pares
sobre una tarea binaria con etiquetas \{apto, no apto\}, construido con fines
puramente ilustrativos ---los valores no provienen de ningún experimento de
esta tesis. La Tabla~\ref{tab:ejemplo-ins} muestra las decisiones de un
modelo hipotético sobre cada miembro del par.

\begin{table}[H]
  \centering
  \caption[Ejemplo ilustrativo del cálculo del \INS de decisión.]{Ejemplo
  ilustrativo (datos ficticios) del cálculo de
  $\mathrm{INS}_{\mathrm{dec}}$.}
  \label{tab:ejemplo-ins}
  \begin{tabular}{clll}
    \toprule
    \textbf{Par} & $f(t_i)$ & $f(t_i')$ & \textbf{¿Coincide?} \\
    \midrule
    1 & apto    & apto    & sí \\
    2 & apto    & no apto & no \\
    3 & no apto & no apto & sí \\
    4 & apto    & apto    & sí \\
    \bottomrule
  \end{tabular}
\end{table}

Tres de los cuatro pares reciben la misma decisión, de modo que
$\mathrm{INS}_{\mathrm{dec}} = 3/4 = 0.75$. El ejemplo permite además
ilustrar qué añade la formulación probabilística: supóngase que en el par~1
el modelo asigna a la clase ``apto'' una probabilidad de $0.91$ para $t_1$
pero de $0.62$ para $t_1'$. La decisión no cambia ---ambas superan el umbral
de $0.5$--- y el par cuenta como consistente para
$\mathrm{INS}_{\mathrm{dec}}$; sin embargo, la divergencia de
Jensen--Shannon entre $(0.91, 0.09)$ y $(0.62, 0.38)$ es de
aproximadamente $0.089$, un desplazamiento de confianza que
$\mathrm{INS}_{\mathrm{prob}}$ sí registra. En tareas donde la probabilidad
alimenta decisiones posteriores ---ordenar candidatos en lugar de solo
aceptarlos o rechazarlos--- ese desplazamiento tiene consecuencias aunque la
etiqueta no cambie.

\section{Aplicación a los resultados experimentales}
\label{sec:ins-aplicacion}

[Reporte de los valores de \INS calculados sobre los modelos del
Capítulo~\ref{cap:resultados}, en sus formulaciones agregada y desagregada
por tipo de variación, con su correspondiente interpretación.]

\section{Interpretación y alcances del indicador}
\label{sec:ins-interpretacion}

Un indicador solo es útil si se conocen con precisión sus condiciones de
lectura. Esta sección delimita qué afirma y qué no afirma un valor del \INS,
bajo qué condiciones es comparable y cómo se sitúa frente a las métricas de
equidad algorítmica existentes.

\subsection{Qué mide y qué no mide el \INS}

El \INS mide \emph{consistencia intra-par}: el grado en que un clasificador
responde igual ante dos textos cuyo contenido es equivalente y cuyo registro
difiere. Cuatro precisiones acotan esta lectura.

Primera, consistencia no es corrección. Ambas dimensiones son ortogonales y
su cruce genera cuatro escenarios cualitativamente distintos: un modelo
exacto y consistente (el caso deseable); un modelo exacto en el agregado pero
inconsistente entre registros, que es precisamente el sesgo que las métricas
agregadas encubren; un modelo inexacto pero consistente, que falla por igual
para todos los registros; y un modelo inexacto e inconsistente. El \INS
distingue las columnas de este cruce, no las filas: por eso debe reportarse
siempre junto a las métricas de desempeño de la tarea, nunca en su lugar.

Segunda, y como caso límite de lo anterior: la neutralidad perfecta es
trivialmente alcanzable. Un clasificador constante, que ignora por completo
su entrada, obtiene $\mathrm{INS}_{\mathrm{dec}} = 1$ por construcción. Un
\INS alto, aislado de toda métrica de desempeño, no acredita nada: el índice
es informativo únicamente para modelos que resuelven la tarea con un
desempeño no trivial.

Tercera, el \INS es simétrico y, por tanto, no direccional: registra
\emph{que} el modelo cambia su respuesta ante el registro, pero no
\emph{a quién} favorece el cambio. Un mismo valor de \INS es compatible con
un modelo que penaliza sistemáticamente el registro coloquial y con uno que
lo favorece. Establecer la dirección del sesgo exige un análisis
complementario ---por ejemplo, comparar la tasa de asignación de la clase
favorable por condición de registro---, que en esta tesis se reporta en la
Sección~\ref{sec:resultados-experimentos}.

Cuarta, el \INS es una medida relativa al \emph{benchmark}, no una propiedad
intrínseca del modelo. Su validez hereda la validez de la operacionalización
de ``equivalencia de contenido'' descrita en la
Sección~\ref{sec:benchmark-diseno}: si los pares confunden variación de
registro con variación de contenido, el índice atribuirá al sesgo
sociolectal inconsistencias que en realidad responden a información
distinta. Todo valor de \INS debe leerse, en rigor, como
``\INS del modelo $f$ \emph{sobre el benchmark} $\mathcal{B}$''.

Finalmente, conviene precisar la relación entre las dos formulaciones. La
probabilística es más sensible: detecta desplazamientos de confianza que no
alcanzan a cruzar el umbral de decisión, como ilustró el ejemplo de la
Sección~\ref{sec:ins-procedimiento}. La de decisión, en cambio, refleja la
consecuencia material: lo que la persona evaluada experimenta es la decisión,
no la distribución interna del modelo. Reportar ambas, cuando es posible,
ofrece la lectura más completa: la probabilística como instrumento de
diagnóstico fino, la de decisión como medida del efecto práctico.

\subsection{Umbrales y comparabilidad}

No existe un umbral universal a partir del cual un valor de \INS certifique
neutralidad, y esta tesis deliberadamente no propone uno. Lo que sí pueden
ofrecerse son puntos de referencia que delimitan el rango interpretable. En
una tarea binaria, un clasificador que respondiera al azar de manera
independiente sobre cada texto obtendría un
$\mathrm{INS}_{\mathrm{dec}}$ esperado de $0.5$; el clasificador constante,
como se señaló, alcanza $1$ de forma trivial. Los valores informativos de un
modelo real se leen contra ese telón de fondo y, sobre todo, en
\emph{comparación}: entre modelos evaluados sobre el mismo \emph{benchmark},
entre versiones sucesivas de un mismo modelo, o entre tipos de variación
dentro de un mismo modelo.

De ahí se siguen las condiciones de comparabilidad. Dos valores de \INS son
directamente comparables solo si provienen del mismo \emph{benchmark} ---los
mismos pares, la misma tarea, el mismo espacio de etiquetas--- y de la misma
formulación (decisión con decisión, probabilística con probabilística).
Comparar valores obtenidos sobre \emph{benchmarks} distintos no es válido sin
más: un conjunto de pares con variaciones sutiles producirá índices
sistemáticamente más altos que uno con variaciones marcadas, sin que ello
diga nada sobre los modelos.

Por último, el perfil desagregado
$\{\mathrm{INS}(f, \mathcal{B} \mid v)\}_{v \in V}$ es más informativo que el
escalar global, pues localiza el problema: un modelo puede ser robusto ante
anglicismos e inconsistente ante variación ortográfica, y esa estructura
---invisible en el agregado--- es la que orienta tanto el diagnóstico como
una eventual mitigación. El protocolo de reporte adoptado en esta tesis
comprende, por ello, el \INS global en sus dos formulaciones cuando es
posible, el perfil desagregado por tipo de variación con el tamaño $N_v$ de
cada celda, y las métricas de desempeño del modelo sobre la tarea.

\subsection{Relación con otras métricas de equidad}

La literatura sobre equidad algorítmica ha producido una familia amplia de
criterios formales, sistematizada en revisiones como la de Mehrabi y
colaboradores \cite{Mehrabi2021Survey}. Situar el \INS en ese mapa ayuda a
precisar su aporte.

Las métricas de \emph{equidad de grupo} ---la paridad demográfica, que exige
tasas de asignación de la clase favorable iguales entre grupos, o la igualdad
de oportunidades, que exige tasas de verdaderos positivos iguales
condicionadas a la etiqueta real--- comparan agregados calculados sobre
individuos \emph{distintos} pertenecientes a grupos distintos
\cite{Mehrabi2021Survey}. El \INS opera en otra unidad de análisis: compara
las respuestas del modelo ante versiones del \emph{mismo} contenido, lo que
lo aproxima a la familia de la \emph{equidad individual}, cuyo principio es
que entradas similares reciban tratamientos similares. El \emph{benchmark} de
pares controlados es, en efecto, una operacionalización textual de
``similitud'': dos textos se declaran similares cuando su contenido es
equivalente y solo difiere su registro.

Esta construcción acerca el \INS, en particular, a la noción de
\emph{equidad contrafáctica}, según la cual una decisión es justa si no
cambiaría en el mundo contrafáctico donde el individuo perteneciera a otro
grupo protegido [FUENTE: añadir la referencia primaria de equidad
contrafáctica (Kusner et al., 2017, \emph{Counterfactual Fairness}, NeurIPS)
a referencias.bib; mientras tanto, la noción está cubierta por el survey de
Mehrabi]. Cada par $(t, t')$ del \emph{benchmark} funciona como un
contrafáctico textual: el ``individuo'' es el contenido informativo y la
intervención consiste en cambiar únicamente el atributo sensible, el
registro. La diferencia importante es de estatus epistémico: la equidad
contrafáctica plantea la intervención sobre un modelo causal del mundo,
mientras que el \INS la realiza materialmente sobre el texto, que es donde el
atributo sociolectal se manifiesta. Esa materialización es lo que vuelve al
índice computable sin supuestos causales adicionales.

El \INS comparte también una afinidad metodológica con las pruebas de
invariancia usadas en la evaluación conductual de modelos de \PLN, en las que
se verifica que perturbaciones que no alteran el significado no alteren la
predicción [FUENTE: añadir la referencia de evaluación conductual por
invariancia (Ribeiro et al., 2020, \emph{CheckList}, ACL) a
referencias.bib]. Bajo esa lectura, el \INS es una prueba de invariancia cuya
perturbación no es sintética ni arbitraria, sino sociolingüísticamente
anclada: las variaciones provienen de los marcadores de registro
documentados en el Capítulo~\ref{cap:marco} y derivados empíricamente de la
encuesta.

Dos ventajas prácticas se siguen de esta posición en el mapa. Primera, a
diferencia de la igualdad de oportunidades, el \INS en su formulación de
consistencia no requiere etiquetas de verdad de terreno: basta el
\emph{benchmark} y el acceso al modelo, lo que lo hace aplicable a sistemas
de caja negra cuya tarea ni siquiera tiene una etiqueta ``correcta''
verificable. Segunda, su unidad de análisis pareada lo protege de un modo de
fallo conocido de las métricas de grupo: un sistema puede satisfacer la
paridad demográfica compensando errores entre subpoblaciones sin tratar de
manera consistente a ningún individuo en particular. Nada de esto convierte
al \INS en sustituto de aquellas métricas: mide una dimensión ---la
consistencia ante variación sociolectal--- que las métricas de grupo no
capturan, y omite otras ---las tasas de acierto condicionadas a la verdad de
terreno--- que ellas sí. Su lugar es el de un complemento dentro de una
batería de evaluación, no el de una medida única de equidad.
